\section{Semantic transparency}\label{sec:4}

Section 3 set up form-meaning transparency as a theoretical gradient and a network property. Section 4 narrows $R$ to interpretation and turns to a quantitative case: rated transparency of English noun-noun compounds. §\ref{sec:4-3} places referential transparency as the schema's boundary case, where $X$ is a context rather than a form.

\subsection{Compound transparency: the Bell and Schäfer model}\label{sec:4-1}

Bell and Schäfer \autocite*{bell-schafer-2016-modelling-transparency} model the perceived semantic transparency of English noun-noun compounds as a function of \mention{expectedness} in their internal semantic structure. A compound is rated as more transparent when its semantic structure is more expected, both in the predictability of its constituents in their positions and in the predictability of the semantic relation between the constituents given those constituents. Transparency is treated as a scalar variable rather than a binary or four-fold property. The empirical anchor is the database of Reddy and colleagues (2011): 30 ratings on a 0--5 scale for each of 90 two-part NN compounds, totalling 8,100 ratings (Appendix B.5).

Bell and Schäfer's models identify three predictors of higher rated transparency. First, \mention{N1 frequency}: compounds whose first noun occurs more frequently in the language overall are rated as more transparent. Second, \mention{semantic-relation expectedness given N1}: compounds whose semantic relation is more frequently associated with that particular N1 are rated as more transparent. Third, \mention{N2 productivity}: compounds whose second noun occurs more often as the head of NN compounds are rated as more transparent. The first and third predictors are corpus-statistical; the second requires hand-coding of semantic relations using Levi's (1978) nine-predicate inventory (CAUSE, HAVE, MAKE, USE, BE, IN, FOR, FROM, ABOUT; Appendix B.1) and aggregation across compounds sharing each N1.

The two headline conclusions: N1 and N2 contribute asymmetrically to whole-compound transparency, and all significant predictors of rated transparency also independently predict processing speed. The latter motivates Bell and Schäfer's tentative hypothesis that ``perceived transparency may itself be a reflex of ease of processing'' \autocite[157]{bell-schafer-2016-modelling-transparency}. The hypothesis is offered as a direction for further work, not a result. Stress-pattern speculation about the IN/FOR positive associations is in Appendix B.3; the N1/N2 asymmetry detail is in Appendix B.4.

\subsection{Compound transparency in the umbrella schema}\label{sec:4-2}

The umbrella schema's variables map directly. $X$ (the mediator) is the compound, e.g.~\mention{car-wash}, \mention{jailbird}, \mention{hogwash}. $P$ (what remains accessible) is the constituent meanings and their semantic relation. $R$ (the relation) is interpretation: rated transparency operationalises how successfully the constituents' meanings reach interpretation through the compound. The schema's ``remains accessible'' requirement is gradient here, just as in §\ref{sec:3}: a compound's parts are accessible to varying degrees, with rated transparency as the empirical proxy.

Bell and Schäfer's three predictors fit naturally inside the schema. N1 frequency, N2 productivity, and semantic-relation expectedness given N1 are properties of the network in which a compound sits, not properties of the compound in isolation. A compound is rated more transparent when its constituents and the relation between them are more conventional in the language. This is the network-relative view §\ref{sec:3-3} framed theoretically, expressed quantitatively. Compound transparency isn't a property of the compound alone but of the compound's position in the constructicon-and-corpus.

Bell and Schäfer's predictors are corpus-statistical; the constructicon network is a theoretical posit. Frequency-driven entrenchment is the standard usage-based story for how the constructicon is shaped, but Bell and Schäfer don't directly measure constructicon-network properties. Their work is consistent with §\ref{sec:3}'s network claim without testing it.

\subsection{Referential transparency as boundary case}\label{sec:4-3}

The schema's referential boundary is independent of the form-meaning accessibility family. \mention{Lois Lane believes Superman is brave} is a textbook intensional context; substituting \mention{Clark Kent} for \mention{Superman} yields a sentence whose truth-conditions may not be preserved (the de re reading allows substitution; the de dicto reading doesn't). This is referential transparency in the philosophical-semantics tradition, and the schema technically applies: $X$ is the intensional context (the embedding under \mention{believe}), $P$ is the referent of the embedded singular term, $R$ is substitution under co-reference.

What makes this the boundary case is that $X$ is a \mention{context} rather than a \mention{linguistic form}. In all the cases §§\ref{sec:2}--\ref{sec:3} worked through, $X$ was a form whose internal structure could have mediated access to $P$: a quantificational-noun-of-complement construction; a compound; a partly-schematic argument-structure construction. An intensional context isn't a form in that sense. It's a configuration that creates substitution-opacity, and the de re/de dicto distinction is about whether substitution is allowed, not about whether form-internal compositional access is preserved.

The schema therefore stretches thinnest at this boundary. $P$ is a referent, not a constituent meaning or a number feature; $R$ is substitution-salva-veritate, not agreement or interpretation. The umbrella applies on a sufficiently abstract reading: referential transparency is a kind of mediated accessibility, where the mediator is the intensional embedding and the referent is what does or doesn't remain accessible. But the difference in the type of $X$ is real and matters for the metaphysical conclusion. Referential transparency isn't part of the form-meaning accessibility family of §§\ref{sec:3}--\ref{sec:5}.

The philosophical literature on referential transparency (Quine \autocite*{quine-1956-quantifiers-propositional-attitudes} as canonical anchor; the Frege/Kripke/Kaplan tradition) would dominate the section if developed. The boundary marks a real heterogeneity in the umbrella, not a defect: the schema's $R$-relativity handles the case without forcing it into the form-meaning family.

\subsection{Compound transparency as a within-domain projectibility profile}\label{sec:4-4}

Compound transparency joins constructional compositionality (§\ref{sec:3}) as a within-domain projectibility profile on form-meaning accessibility. Class-membership talk doesn't apply at this level: there is no small, conventionally stable class of ``transparent compounds'' in the way there is for number-transparent quantificational nouns. The story here is gradient: compounds are more or less transparent on a continuous scale, and the predictors of higher transparency are corpus-statistical and constructional.

Compound transparency licenses non-trivial within-domain projection. Knowing a compound's N1 frequency, semantic-relation expectedness, and N2 productivity lets us predict its rated transparency, with non-trivial fit in Bell and Schäfer's models. The stabilising source is usage-based entrenchment: the network of compounds and their semantic relations is shaped by frequency, productivity, and the kind of conventional patterns that emerge from repeated use. Falsification condition: if N1 frequency, semantic-relation expectedness, and N2 productivity failed to predict rated transparency in an extension or replication of Bell and Schäfer's models, the within-domain profile would fail. Whether the within-domain projection joins with §\ref{sec:3}'s theoretical profile and §\ref{sec:5}'s processing window into a cross-domain projectibility cluster on shared items is left to §\ref{sec:7}, and §\ref{sec:7} doesn't take that result as delivered.

What §\ref{sec:4} contributes that §\ref{sec:3} didn't is empirical content. The compositionality gradient theorised in §\ref{sec:3} is a measurable property whose predictors are corpus-derivable and whose values correlate with rated transparency. That correlation is what makes compound transparency a quantitative window onto the gradient. Whether the same predictors would also track processing latencies, which would strengthen the case for Khalidi-style causal-hierarchical refinement (without by itself establishing it), is a question for §5.

\subsection{Bridge to §\ref{sec:5}}\label{sec:4-5}

§\ref{sec:5} takes the third window on the family: time-sensitive lexical processing. Where Bell and Schäfer measure rated transparency, Heyer and Kornishova measure masked priming. Where rated transparency is a steady-state property, priming is a time-course property. Whether the windows of §\ref{sec:3}, §\ref{sec:4}, and §\ref{sec:5} jointly support the kind of causal-hierarchical projection that Khalidi's framework formalises is the question §\ref{sec:7} will return to.
